% =============================================================================
\chapter{Conclusion}
\label{ch:conclusion}
% =============================================================================

This thesis evaluated three mobile pose estimation models and six additional variants on a rehabilitation-oriented proxy benchmark to study how well standard benchmark rankings carry into rehabilitation-oriented use. The evaluation covered body-joint accuracy, temporal stability, viewpoint sensitivity, multi-person behavior, detection completeness, inference speed, and cross-dataset transferability.

Two findings stand out. First, on REHAB24-6, MoveNet MultiPose is the strongest main-model default for body-joint analysis: it achieves the lowest subject-cluster error and the lowest frame-to-frame prediction-displacement value. Second, the three models do not collapse into a single ranking. YOLOv8 is the most detection-complete but also the most exposed to downstream disambiguation, MoveNet is the most conservative in secondary-person exposure yet does not show the lowest coach-scene failure rate, and MediaPipe shows the lowest catastrophic-failure rate on the coach extreme-case subset despite lower body-joint accuracy. The data therefore support three architectural profiles rather than a universal winner, and secondary-person detection frequency alone does not explain coach-scene robustness. A failure-mode decomposition reinforces this: each main model is dominated by a distinct failure category consistent with its architectural paradigm (Keypoint-Displacement for MediaPipe, Confidence-Collapse for MoveNet, Missing-Detection for YOLOv8).

The results should still be read within three limits. REHAB24-6 contains ten healthy adult volunteers performing physiotherapist-guided rehabilitation exercises rather than a clinical patient cohort. The cross-dataset comparison on COCO val2017 uses a different evaluation protocol, so the ranking shift should be read as limited transferability rather than as evidence of benchmark invalidity. All models were evaluated frame-independently; the reported person-switch and temporal-stability values therefore describe unsmoothed model outputs rather than a fully engineered deployment pipeline.

For CPU-based body-joint rehabilitation analysis, MoveNet MultiPose is the strongest default among the three main models. Other configurations are preferable when different criteria dominate: YOLOv8 for detection completeness, MediaPipe for lower coach-scene failure rate, richer 33-landmark output, or simpler integration, and MoveNet\_SP\_Lightning for maximum throughput. Moving from this proxy benchmark toward clinical rehabilitation use will require validation on real patient data and cleaner cross-dataset protocols.
